\section{Experiments and Results}
\label{sec:experiments}

The experiment has three evidence layers: offline data and simulator
acceptance, adaptive development, and formal confirmation. Only the last layer
is used for the confirmatory effect claim.

\subsection{Target benchmark and offline data}
\label{sec:exp:benchmark}

The target benchmark uses SUMO on the 12 directional services in
\Cref{fig:network_topology}. Each 24,000-s episode contains the same timetable,
vehicle capacities, road network, and 389 scheduled vehicle trips. Passenger
demand is evaluated at multipliers 0.75, 1.00, and 1.25. Within an environment
block, the three demand levels share one demand-realization seed and one SUMO
seed; departure times receive stable-uniform perturbations of at most 120 s.

The immutable offline file contains 3,284,718 chronological rows from 240
complete episodes. Ten independent environment blocks crossed with three
demand levels and eight behavior policies give 30 exogenous cells and 240
trajectories. Behaviors are zero hold, two random hurdle policies, four
fixed-gain schedule-adjusted headway policies, and a closed-form myopic-welfare
policy. They provide action coverage and are not labelled experts. Data
collection used SUMO 1.27.1 and Python 3.10.12.

The calibrated online simulator shares the timetable, capacity, journey,
state, action, and reward contracts. In a preregistered zero-policy audit over
30 audited cells, every 95\% bootstrap interval stayed within its threshold.
The registered relative tolerances were 3\% for passenger demand, 5\% for
decision count, 15\% for generalized and waiting time, 10\% for in-vehicle
time, and 5\% for completed vehicle-trip duration; completion used an absolute
0.03 tolerance.
The largest point gaps were 4.61\% for generalized passenger time, 8.75\% for
waiting time, 1.66\% for in-vehicle time, 0.43 percentage points for
completion, and 1.04\% for completed vehicle-trip duration. A second audit
crossed the same cells with seven nonzero interventions. The largest absolute
95\% endpoint was 0.00574 for completion, 0.0390 for normalized generalized
passenger-time response, 0.0498 for waiting response, 0.0200 for in-vehicle
response, 0.2078 for relative mean hold, and 0.0117 for vehicle-trip duration;
all were below their preregistered limits. Intervention-response tolerances
were 0.03 for absolute completion response, 0.10 for normalized generalized,
waiting, and in-vehicle responses, and 0.05 for decision and trip-duration
responses. Policy-level mean hold and exact-zero-action rate used relative
0.35 and absolute 0.10 limits, respectively. Both audits used 10,000 bootstrap
resamples of episode-level quantities within each demand level, resampling
SUMO and simulator samples separately. Intervention outcomes were first
differenced against the same environment's zero-hold outcome; time and count
responses were normalized by the corresponding SUMO zero-hold mean. These
acceptance audits precede the independent formal test below.

\subsection{Formal confirmation matrix}
\label{sec:exp:formal_design}

Formal evaluation was endpoint-only: no area under a learning curve (AULC) was
computed or claimed. The sole candidate was the reserve-augmented policy in
\Cref{eq:completion_reserve}. Its direct behavior parent was
\Cref{eq:direct_parent} with the final 60-s reserve disabled. Five
seed-matched actor tensors were fixed before evaluation.

Seven supportive controls were also frozen: the raw v14 terminal-quarter
actor, the base H2O+ endpoint, online-only SAC, delayed-actor v10, actor-hold
v11, actor-ramp v12, and proximal-strong v13. Together with the candidate and
direct parent, nine algorithm labels crossed with five training seeds to form
45 physical behaviors. These controls provide lineage context; they do not
expand the confirmatory claim beyond candidate versus direct parent.

Ten fresh formal environment blocks used SUMO seeds 8101--8110 and demand
realization seeds 8201--8210. None appeared in training or adaptive
development. Crossing 10 blocks, three demand levels, and 45 behaviors produced
1,350 new SUMO episodes and 30 fresh environment cells. Formal execution used
SUMO/libsumo 1.25.0 under Python 3.8.20. All policies in a cell shared the same
exogenous realization.

\subsection{Endpoints and preregistered inference}
\label{sec:exp:statistics}

The primary endpoint is $Y$, accrued generalized passenger seconds per departed
passenger; lower is better. The service endpoint is $R_{\mathrm{cmp}}$, and
unfinished passenger count $U$ is reported descriptively, all as defined in
\Cref{eq:completion_rate}. The independent analysis unit is the formal environment block.
For each candidate--parent comparison, the five training seeds and three demand
levels yield 15 paired observations per block; these are averaged before any
sign test. Let $\beta\in\{1,\ldots,10\}$ index formal blocks,
$\xi\in\{1,\ldots,5\}$ index training seeds, and
$d\in\mathcal{D}_{\mathrm{dem}}=\{0.75,1.00,1.25\}$ denote the demand
multiplier. Superscripts $\mathrm{safe}$ and $\mathrm{par}$ identify the
candidate and parent, respectively. Their primary block difference is
\begin{equation}
\Delta Y_{\beta}=\frac{1}{15}
\sum_{\xi=1}^{5}\sum_{d\in\mathcal{D}_{\mathrm{dem}}}
\left(Y^{\mathrm{safe}}_{\beta,\xi,d}
-Y^{\mathrm{par}}_{\beta,\xi,d}\right).
\label{eq:block_difference}
\end{equation}
The same averaging in \Cref{eq:block_difference} defines
$\Delta R_{\mathrm{cmp},\beta}$ and $\Delta U_{\beta}$. In the tables,
$\Delta Y$, $\Delta R_{\mathrm{cmp}}$, and $\Delta U$ denote mean paired
differences over all cells or the stated demand or seed subset.

Primary superiority uses a one-sided exact paired sign test. Completion uses
the same test after shifting each block difference by the noninferiority margin
$\varepsilon_{\mathrm{NI}}=0.003$: a block is a win when
$\Delta R_{\mathrm{cmp},\beta}>-\varepsilon_{\mathrm{NI}}$.
A primary win requires $\Delta Y_{\beta}<0$. Ties are nonwins. For a win
count $w\in\{0,\ldots,10\}$, the one-sided sign-test probability is
\begin{equation}
p_{\mathrm{sign}}(w)=2^{-10}\sum_{\upsilon=w}^{10}
\binom{10}{\upsilon},
\label{eq:sign_probability}
\end{equation}
where $\upsilon$ is the summation index for possible win counts and
$\binom{10}{\upsilon}$ is the binomial coefficient. At
one-sided significance level $\alpha_{\mathrm{sig}}=0.025$, at least 9 of
10 wins are required by \Cref{eq:sign_probability}
($p_{\mathrm{sign}}(9)=0.010742$). The confirmatory intersection--union test passes
only if both primary superiority and completion noninferiority pass. Mean
direction, demand-stratum direction, result completeness, and parent artifact
and training stability were additional preregistered gates.

\subsection{Confirmed candidate--parent effect}
\label{sec:exp:primary_results}

All 1,350 results and all ten 135-job block manifests passed independent
validation. \Cref{tab:formal_primary} reports the endpoint means. The candidate
reduced the primary endpoint by 92.12 s per departed passenger, a descriptive
relative reduction of 3.71\%. Completion increased by 0.00466 (0.466
percentage points), and unfinished passengers decreased by 77.22 (11.78\%).

\begin{table}[t]
\centering
\small
\caption{Formal endpoint means for the preregistered candidate and its direct
behavior parent. The difference is candidate minus parent.}
\label{tab:formal_primary}
\resizebox{\linewidth}{!}{%
\begin{tabular}{lrrr}
\toprule
Policy & Gen. passenger s/departed & Completion & Unfinished \\
\midrule
Direct parent & 2,484.86 & 0.96185 & 655.42 \\
Reserve-augmented candidate & \textbf{2,392.73} & \textbf{0.96650} & \textbf{578.20} \\
Difference & $-92.12$ & $+0.00466$ & $-77.22$ \\
\bottomrule
\end{tabular}
}
\end{table}

The primary difference was negative in all ten blocks, giving 10 of 10 wins
and a one-sided exact sign-test value of $p_{\mathrm{sign}}(10)=0.0009766$. Completion differences
were positive in every block, so completion noninferiority also achieved 10 of
10 wins and $p_{\mathrm{sign}}(10)=0.0009766$ after the margin shift. The intersection--union test
therefore passed. All preregistered artifact, parent-stability, mean,
demand-stratum, and result-completeness gates passed as well.

The effect was directionally consistent across demand levels
(\Cref{tab:demand_effect}). Its magnitude increased with demand: the primary
reduction ranged from 40.29 s at demand 0.75 to 148.46 s at demand 1.25, while
completion improved at every level. It was also directionally consistent for
each of the five training seeds: primary differences ranged from $-87.25$ to
$-97.44$ s and completion differences from $+0.00444$ to $+0.00485$
(\Cref{app:seed_effect,tab:seed_effect}).

\begin{table}[t]
\centering
\small
\caption{Candidate-minus-parent formal differences by demand multiplier.
Negative generalized time and unfinished counts are favorable; positive
completion is favorable.}
\label{tab:demand_effect}
\begin{tabular}{lrrr}
\toprule
Demand & $\Delta Y$ (s/departed) & $\Delta R_{\mathrm{cmp}}$ & $\Delta U$ (passengers) \\
\midrule
0.75 & $-40.29$ & $+0.00185$ & $-21.00$ \\
1.00 & $-87.61$ & $+0.00470$ & $-70.92$ \\
1.25 & $-148.46$ & $+0.00741$ & $-139.74$ \\
\bottomrule
\end{tabular}
\end{table}

\subsection{Supportive lineage controls}
\label{sec:exp:controls}

\Cref{tab:formal_controls} gives descriptive means for all fixed algorithm
labels. The reserve-augmented candidate had lower generalized passenger time and
higher completion than every listed control, including online-only SAC. Its
primary block difference versus each control was favorable in all ten blocks.
These comparisons were prespecified robustness gates, but the controls do not
share the candidate's exact behavior parent and are not part of the
intersection--union confirmatory claim.

\begin{table*}[t]
\centering
\small
\caption{Descriptive formal means for the fixed candidate, direct parent, raw
actor, and lineage controls. Bold denotes the best displayed mean.}
\label{tab:formal_controls}
\resizebox{\textwidth}{!}{%
\begin{tabular}{lrrr}
\toprule
Algorithm label & Gen. passenger s/departed & Completion & Unfinished passengers \\
\midrule
Reserve-augmented candidate & \textbf{2,392.73} & \textbf{0.96650} & \textbf{578.20} \\
Direct behavior parent & 2,484.86 & 0.96185 & 655.42 \\
Raw v14 terminal-quarter actor & 2,623.19 & 0.96453 & 610.03 \\
Base H2O+ & 2,696.98 & 0.96288 & 637.41 \\
Online-only SAC & 2,593.85 & 0.96232 & 648.31 \\
Delayed actor (v10) & 2,532.20 & 0.95894 & 702.57 \\
Actor hold (v11) & 2,569.68 & 0.96023 & 681.31 \\
Actor ramp (v12) & 2,585.56 & 0.96194 & 653.26 \\
Proximal strong (v13) & 2,581.26 & 0.96053 & 676.66 \\
\bottomrule
\end{tabular}
}
\end{table*}

\subsection{Data and environment diagnostics}
\label{sec:exp:validity}

The artifact checks found no nonfinite rows, no passenger-ledger reconstruction
error, and 240 terminal rows for 240 trajectories. Nevertheless, four factors
qualify algorithmic interpretation. \Cref{app:support} provides the support
counts and dataset-integrity details.

First, the 3.28 million serialized transitions arise from only 30 independent
exogenous cells, each reused by eight behavior policies. Statistical evidence
must therefore be organized by environment block rather than transition.
Second, 80.2246\% of offline actions are exact zero. The final low-load guard
appears in 3.91\% of offline rows, and only 3.69\% of those guard-state actions
are at least 30 s. The 60-s reserve should consequently be interpreted as a
bounded operational overlay, not a behavior learned from dense action support.

Third, the ratio diagnostic generalized poorly. Mean joint classifier area
under the receiver-operating-characteristic curve (AUC) was 0.9477 on training data
and 0.4520 on trajectory-held-out validation data; mean raw validation
effective sample size (ESS) fraction was 0.0615. \Cref{app:ratio,tab:ratio_diag}
reports the across-seed diagnostic summary. Unit simulator weights prevented
this failure from changing the sealed optimizer, but no dynamics-correction
effect is established. Fourth, offline collection used SUMO 1.27.1 whereas
formal evaluation used SUMO/libsumo 1.25.0. The fresh formal result therefore
remains favorable in the evaluation runtime, but a paired version-sensitivity
experiment was not conducted. Canonical feature clipping was
otherwise limited: only 1,565 rows (0.0476\%) exceeded the backward-headway
cap, and the other audited caps had zero exposure.

\subsection{Execution and reproducibility}
\label{sec:exp:reproducibility}

Each formal result binds the method identity, training seed, actor tensor,
checkpoint, target dataset, evaluator source, scenario, and runtime. Ten matrix
manifests account for all 1,350 files. A frozen analysis produced the reported
statistics, and a separate implementation independently recomputed the result
count, means, block wins, sign-test value, and all preregistered gates. The
closeout also verified that no local SUMO execution occurred during analysis
and that formal result files were not modified.

Only result JavaScript Object Notation (JSON) files, matrix manifests, compact logs, and verification
metadata were synchronized from the evaluation server. Checkpoints, datasets,
optimizer state, replay buffers, and remote workspaces were not duplicated in
the paper artifact. This preserves a minimal reproducibility package while the
frozen deployment manifest records the larger server-side dependencies.
\Cref{app:evidence} identifies the compact analysis and verification files.
